\documentclass[aspectratio=169,11pt]{ctexbeamer}

\usetheme{Madrid}
\usecolortheme{default}
\usefonttheme{professionalfonts}
\setbeamertemplate{navigation symbols}{}
\setbeamertemplate{footline}[frame number]

\usepackage{booktabs}
\pdfstringdefDisableCommands{\def\texttt#1{#1}}

\definecolor{HustBlue}{RGB}{17,64,124}
\definecolor{HustLight}{RGB}{233,241,250}

\setbeamercolor{palette primary}{bg=HustBlue,fg=white}
\setbeamercolor{palette secondary}{bg=HustBlue,fg=white}
\setbeamercolor{palette tertiary}{bg=HustBlue,fg=white}
\setbeamercolor{structure}{fg=HustBlue}
\setbeamercolor{frametitle}{bg=HustLight,fg=HustBlue}
\setbeamercolor{title}{fg=HustBlue}
\setbeamercolor{block title}{bg=HustBlue,fg=white}
\setbeamercolor{block body}{bg=HustLight,fg=black}

\title{课题组阶段工作与下一步安排}
\subtitle{团队会议发言材料}
\author{张书豪}
\date{2026-04-01}

\begin{document}

\begin{frame}
  \titlepage
\end{frame}

\begin{frame}{先把总判断统一下来}
  \begin{itemize}
    \item 总体实施方案和技术架构已经完成阶段性定型。
    \item \texttt{vllm-hust} 已经具备统一推理底座的基础能力。
    \item \texttt{vllm-hust-workstation} 已经是应用平台，不是简单入口。
    \item \texttt{EvoScientist} 已经作为首个真实多智能体业务对象接入本地链路。
    \item 当前重点已经从“能不能跑”转到“场景验证、数据闭环和跨硬件实测”。
  \end{itemize}

  \vspace{0.5em}
  \begin{block}{统一判断}
    现在底座已经有了，下一阶段重点不是再讲新架构，而是把真实场景、统一指标和验证证据补齐。
  \end{block}
\end{frame}

\begin{frame}{第一个里程碑怎么界定}
  \begin{itemize}
    \item 由课题组主承担。
    \item 不等于 9 个优化点全面展开。
    \item 核心是公共基线、统一指标、统一评测、阶段材料归口。
    \item 9 个优化点原有分工不变。
  \end{itemize}

  \vspace{0.5em}
  \begin{block}{统一表述}
    第一个里程碑先由课题组把公共底座和统一口径收住，各位老师继续按既定优化点推进。我们不是收回优化点，而是先把后续共用的基线、指标和材料体系搭起来。
  \end{block}
\end{frame}

\begin{frame}{现在已经能明确讲出去的内容}
  \begin{itemize}
    \item 仓库和平台分工已经清楚：\texttt{vllm-hust} 是推理引擎，\texttt{vllm-hust-workstation} 是应用平台，\texttt{vllm-hust-benchmark} 负责评测分析，\texttt{vllm-ascend-hust} 负责国产算力适配，\texttt{vllm-hust-docs} 负责文档和证据归档。
    \item 底座能力已经具备：多模态、流式输出、reasoning、工具调用、结构化输出、benchmark 与 metrics 入口都已经有了。
    \item 真实对象已经接进来了：\texttt{EvoScientist} 已经可以作为首个真实多智能体 workload，默认桥接链路也已经打通。
    \item 画像和国产算力分析都已经有抓手：TTFT、ITL、TPOT、token throughput、error type、reasoning length、tool density 已可开始做事。
  \end{itemize}
\end{frame}

\begin{frame}{边界也要一次讲清楚}
  \begin{itemize}
    \item 我们已经有真实多智能体业务对象，但还不是通用多智能体仿真平台。
    \item 我们已经有 serving 侧数据面和 rollout 方向，但还不是完整流式 RL 训练闭环。
    \item 我们已经有云原生部署和接入基础，但还不是跨集群调度验证完成。
    \item 我们已经识别国产算力主要瓶颈，但还没有把跨硬件同任务实测证据补齐。
  \end{itemize}
\end{frame}

\begin{frame}{课题组下一阶段主任务}
  \begin{itemize}
    \item 把第一个里程碑总收口，统一口径、统一材料、统一数据附件归口。
    \item 把统一指标和统一基线定下来，先把 TTFT、TPOT、吞吐、MFU、单位 token 成本口径收住。
    \item 组织首轮国产硬件基础链路验证，先把“能跑通、有留痕、可归档”的 baseline 做出来。
    \item 把整体架构说明写成正式文本，后面不再反复解释仓库和平台关系。
    \item 形成阶段性答复材料 V1，把已完成、已具备基础、待补证据三层说法固定下来。
  \end{itemize}
\end{frame}

\begin{frame}{其他老师这边重点配合什么}
  \begin{itemize}
    \item 提交各自优化点当前进展、问题和预期指标。
    \item 提供平台环境、样例模型、测试条件和已有数据。
    \item 按统一指标口径补 baseline 或对比结果。
    \item 配合各优化点后续接入统一评测、统一观测和统一材料归档体系。
    \item 提供阶段材料撰写所需的书面输入。
  \end{itemize}

  \vspace{0.5em}
  \begin{block}{配合要求}
    各位老师继续按原优化点分工推进，课题组先把统一口径、统一基线和阶段材料收住；后续按统一模板收结果。
  \end{block}
\end{frame}

\begin{frame}{接下来两周怎么推进}
  \begin{columns}[T]
    \begin{column}{0.32\textwidth}
      \begin{block}{本周}
        \begin{itemize}
          \item 固定材料目录与责任归口
          \item 收齐各优化点状态输入
          \item 固化统一指标与基线脚本
        \end{itemize}
      \end{block}
    \end{column}
    \begin{column}{0.32\textwidth}
      \begin{block}{下周}
        \begin{itemize}
          \item 形成阶段性答复材料 V1
          \item 补首轮国产硬件 baseline
          \item 启动各优化点接口对齐
        \end{itemize}
      \end{block}
    \end{column}
    \begin{column}{0.32\textwidth}
      \begin{block}{两周内}
        \begin{itemize}
          \item 阶段性指标答复材料 V1
          \item 基线数据包
          \item 协同接口清单
        \end{itemize}
      \end{block}
    \end{column}
  \end{columns}
\end{frame}

\begin{frame}{最后收口}
  \begin{block}{会议结论}
    这阶段最大的进展，不是又提出了一套新架构，而是把底层 serving 底座、上层接入入口和首个真实业务对象连起来了。下一步不再泛泛谈架构，而是补真实场景、统一指标、跨硬件实测和数据闭环。第一个里程碑由课题组主承担，先把公共基线和统一材料收住；各位老师按既定优化点继续推进，我们按统一口径来收结果。
  \end{block}
\end{frame}

\end{document}